\subsection*{Interlacing when the smaller part increases}
For $2\le m\le32$ and $n\ge m$, the ordinates of the zeros of
$E_{m,n}$ strictly interlace those of $E_{m+1,n}$. Thus the line
property holds also for every $K_{33,n}$ with $n\ge33$. Consecutive-$n$
interlacing within width $33$ remains outside this certificate.
In particular, the interlacing of $E_{3,n}$ and $E_{4,n}$ follows
unconditionally for $n\ge3$. The cases $n=1,2$ of this particular
interlacing follow from the cross-polytope recurrence and the
larger-part interlacing at width two, respectively.

To see this, keep the pair $Q_n,Q_{n-1}$ fixed, and use subscripts
$(m,n)$ for the short coordinates. In the residue formula take
$C=A_{m+1,n}$ and $D=B_{m+1,n}$. The determinant
\[
 V(n,y^2)=A_{m,n}(y) B_{m+1,n}(y)-B_{m,n}(y) A_{m+1,n}(y)
\]
has even parity. Its degree in $n$ is at most $2r+2$, as the termwise
degree bounds for its factors are $(r,r+2)$ and $(r+1,r+1)$.
The script \texttt{verify\_smaller\_part.py} tests the coordinate bounds
for width $m+1$ by the same interpolation identity; its upper endpoint
includes width $33$. It certifies $V(m+s,z)>0$ for every integer
$s\ge0$ and real $z\ge0$ by applying the Sturm test to the Newton
coefficient polynomials in $s$. The residue sign argument gives the
claimed interlacing. This argument needs the positive Bezoutian only
in the lower width. Run the script with \texttt{2 32}.

\subsection*{A worked example and an independent check from the polytope}
For $K_{6,6}$, dividing both short coordinates by their common factor $132$ gives
\[
 \begin{aligned}
 A(y)&=693y^5-29984y^3+175507y,\\
 B(y)&=19656y^4-442620y^2+881100 .
 \end{aligned}
\]
Thus positivity of a five-by-five real form suffices for the degree-eleven Ehrhart polynomial in this example. Its first four lattice-point counts, at dilations $0,1,2,3$, are $1,73,1915,26547$.

There is a direct way to check these counts. Split a lattice vector into its coordinates $(u,v)$ on the two sides of the graph. Its membership in $kP_{K_{m,n}}$ is equivalent to
\[
       {\textstyle\sum} u_i+{\textstyle\sum} v_j=0,\qquad
       \max(\|u\|_1,\|v\|_1)\le k.
\]
Indeed, a weighted sum of signed edges is a matrix with row sums $u$, column sums $-v$, and cost equal to the sum of the absolute entries. Both marginal norms bound that cost from below. For the converse let $T$ be the larger marginal norm and let $s=\sum u_i$. On each side choose nonnegative positive and negative margins with total positive mass $(T+s)/2$ and total negative mass $(T-s)/2$. Starting from the ordinary positive and negative parts, add the same padding to both where needed. Nonnegative matrices with each pair of prescribed margins exist by successively matching supplies and demands. The difference of the two matrices realizes the desired vector with cost at most $T$. Real edge weights suffice for membership in the polytope.

The independent script \texttt{check\_instances.py} counts integer vectors by their sum and $\ell^1$ norm on each side. It recovers the entire polynomial from the required $d+1$ integer values for $K_{2,2},K_{3,5},K_{6,6},K_{7,9}$, a total of 40 comparisons. The proof for unbounded part sizes uses the Newton and Sturm certificates above.
